\VerdierPhysicalPageAnchor{0029}%
\VerdierSemanticAnchor{intro-7-duality-D4}
\noindent\textit{D4)} When $X$ is a closed subset of $Y$ and $f$ is the
canonical inclusion, the functor $f^!$ is the total derived functor of the
“sections with support in $X$” functor.

\VerdierSemanticAnchor{intro-7-traditional-duality}
This theorem, and its variants in the other contexts indicated above,
contains the traditional formulations of duality: the Poincaré duality
theorem in the topological case, and the Serre duality theorem for coherent
algebraic sheaves on a smooth projective variety. But the use of the language
of derived categories does not merely make it possible to prove, in greater
generality, theorems in which one might have been interested only in simple
special cases. For example, at present one knows how to prove the duality
theorem in its traditional form (of Serre-duality type), for coherent
algebraic sheaves on a smooth proper nonprojective variety over an
algebraically closed field, only by way of the general duality theorem. The
same holds for duality in étale cohomology, for nonsingular projective
varieties.

\VerdierSemanticAnchor{intro-8}
\noindent\textbf{8.}\quad
The duality theorems in cases 1), 2), and 4) lead to nonclassical Lefschetz
fixed-point formulas involving sheaves that may have singularities [10],
[11], [12]. Let $X$ be a compact topological space of finite cohomological
dimension, $A$ a noetherian ring, and $F$ a sheaf of $A$-modules on $X$ with
the following properties:

\VerdierSemanticAnchor{intro-8-lefschetz-hypotheses}
\begin{enumerate}
\renewcommand{\labelenumi}{\alph{enumi})}
\VerdierSemanticAnchor{intro-8-hypothesis-a}
\item At every point of $X$, the stalks of the sheaf $F$ are finitely
generated $A$-modules of finite projective dimension.

\VerdierSemanticAnchor{intro-8-hypothesis-b}
\item A regularity property in a neighborhood of every point of $X$,
analogous to the properties described by Wilder [9].
\end{enumerate}

\VerdierSemanticAnchor{intro-8-endomorphism}
Let $\Phi=(\Phi_X,\Phi_F)\colon(X,F)\longrightarrow(X,F)$ be an endomorphism
of the topological space $X$ equipped with the sheaf $F$. The endomorphism
$\Phi$ induces an endomorphism of a complex of $A$-modules:
\VerdierSemanticAnchor{intro-8-rh-phi}
\[
\mathrm{RH}(\Phi)\colon
\mathrm{RH}(X,F)\longrightarrow\mathrm{RH}(X,F).
\]

\VerdierSemanticAnchor{intro-8-perfect-complex}
Under hypotheses a) and b), the complex $\mathrm{RH}(X,F)$ is a
\emph{perfect complex} of $A$-modules, i.e. it is isomorphic, in the derived
category, to a finite complex whose components are finitely generated
projective $A$-modules.

\VerdierSemanticAnchor{intro-8-lefschetz-number}
One can therefore define the Lefschetz number $\chi(\mathrm{RH}(\Phi))$ as
the alternating sum of the traces of the components of
$\mathrm{RH}(\Phi)$, which are endomorphisms of finitely generated
projective modules. Suppose now, to fix ideas, that the fixed points of
$\Phi$ are isolated. One then has the equality
\VerdierSemanticAnchor{intro-8-lefschetz-formula}
\begingroup
\makeatletter
\tagsleft@true
\let\veqno\@@leqno
\makeatother
\begin{equation}
\chi\bigl(\mathrm{RH}(\Phi)\bigr)
=\sum_{P\ \mathrm{fixed}}\chi_P(\Phi)
\tag{Lef}
\end{equation}
\endgroup
